% main.tex -- CS601T Deep Learning, Programming Assignment 4 (Group 11).
% "Autoencoders for Representation Learning and Image Reconstruction".
% Compile twice with pdflatex -interaction=nonstopmode -halt-on-error main.tex, or latexmk -pdf main.tex.
%
% FIGURES.  Every image goes through \maybeinc (defined below), which embeds the file when it
% exists in the archive root declared by \graphicspath, namely ../results_final/plots/, and
% otherwise draws a labelled placeholder naming the expected path.  The file therefore compiles
% cleanly whether the PNG archive is complete or absent.
% RESULTS.  Every measured number below is taken from results_final/*.json of the final run,
% and each was cross-checked by recomputing it from the archived checkpoint.
%
% Provenance note: this file is derived from main.tex, which remains the template.  main.tex
% carried the measured numbers as \RESULT / \RESULTPCT placeholders and is left that way; the
% filled, corrected version is this file, which compiles to report.pdf.

\documentclass[11pt,a4paper]{article}
\usepackage[utf8]{inputenc}
\usepackage[T1]{fontenc}
\usepackage{lmodern}                 % scalable Type1 fonts (microtype expansion needs them)
\usepackage[margin=1in]{geometry}
\usepackage{amsmath,amssymb,amsthm}
\usepackage{booktabs}
\usepackage{graphicx}
\usepackage[font=small,labelfont=bf,skip=4pt]{caption}
\usepackage{subcaption}
\usepackage{float}
\usepackage{xcolor}
\usepackage{enumitem}
\usepackage{microtype}
\usepackage{array}
\usepackage{longtable}
\usepackage{listings}
\usepackage{fancyhdr}
\usepackage{siunitx}
\usepackage[colorlinks=true,linkcolor=blue!55!black,citecolor=blue!55!black,urlcolor=blue!55!black]{hyperref}

\graphicspath{{../results_final/plots/}}            % the results archive root
\newcommand{\figsrc}[1]{../results_final/plots/#1}
\newcommand{\maybeinc}[2][]{\IfFileExists{\figsrc{#2}}{\includegraphics[#1]{#2}}{\figplaceholder{#2}}}
\newcommand{\figplaceholder}[1]{\fbox{\begin{minipage}[c][2.8em][c]{0.78\linewidth}\centering
  \footnotesize\texttt{\detokenize{#1}}\\[1pt]\scriptsize figure not archived yet\end{minipage}}}
\newcommand{\R}{\mathbb{R}}\newcommand{\given}{\mid}\newcommand{\Acc}{\mathrm{Acc}}
\newcommand{\code}[1]{\texttt{#1}}
% \RESULT / \RESULTPCT are retained so that a future run can be dropped back into the
% template without editing the macro preamble; they are unused here.
\setlength{\parskip}{3pt}\setlength{\parindent}{0pt}\renewcommand{\arraystretch}{1.13}
\lstset{language=Python,basicstyle=\small\ttfamily,keywordstyle=\bfseries,commentstyle=\itshape,
        frame=single,breaklines=true,showstringspaces=false,numbers=left,numberstyle=\tiny,
        numbersep=6pt,xleftmargin=13pt,aboveskip=5pt,belowskip=5pt}
\pagestyle{fancy}\fancyhf{}
\fancyhead[L]{\small\itshape CS601T: Deep Learning --- Assignment 4}
\fancyhead[R]{\small\itshape Group 11}
\fancyfoot[L]{\small\itshape CS601T: Deep Learning --- Assignment 4}
\fancyfoot[C]{\small\itshape\thepage}
\fancyfoot[R]{\small\itshape Group 11}
\renewcommand{\headrulewidth}{0.4pt}\renewcommand{\footrulewidth}{0.4pt}

\begin{document}

\begin{titlepage}
\centering
\vspace*{1em}
{\large\scshape CS601T: Deep Learning\par}\vspace{0.2em}{\normalsize Programming Assignment 4\par}
\vspace{1.2em}\rule{\linewidth}{1.0pt}\\[0.6em]
{\LARGE\bfseries Autoencoders for Representation Learning\\[0.2em] and Image Reconstruction\par}
\vspace{0.6em}\rule{\linewidth}{1.0pt}\vspace{2.4em}
{\large\bfseries Group 11\par}\vspace{0.8em}
\begin{tabular}{@{}ll@{}}
\toprule\textbf{Roll Number} & \textbf{Name} \\ \midrule
CS24BT055 & Himanshu Nagwanshi \\
CS24BT060 & Anshuman Rahul \\
CS24BT061 & Sidharth B \\
\bottomrule
\end{tabular}
\vspace{2.4em}
\begin{tabular}{@{}ll@{}}
\toprule\textbf{Course} & CS601T: Deep Learning \\ \midrule
\textbf{Semester} & August--November 2026 \\
\textbf{Institute} & Indian Institute of Technology Delhi \\
\bottomrule
\end{tabular}
\vfill
{\small Code repository: \code{Group11\_Assignment4\_code} \quad$\vert$\quad Typeset with pdf\LaTeX\par}
\vspace{1em}
\end{titlepage}

\begin{abstract}
\noindent
This report studies autoencoders as a route to supervised representation learning on a five-class
MNIST subset (digits $0,4,5,6,7$, $28\times28$ images flattened to $D=784$).  We first establish a
linear reference point: principal component analysis fitted on the training split only, reducing
$\mathbf{x}$ to $k\in\{32,64,128,256\}$ components, classified by a fully connected network (FCNN).
We then train two autoencoder families, a one-hidden encoder/decoder pair
($784\rightarrow k\rightarrow784$) and a three-hidden family
($784\rightarrow400\rightarrow k\rightarrow400\rightarrow784$), each with a strictly linear
bottleneck of width $k$ and logistic sigmoid elsewhere, minimising the mean squared reconstruction
error with Adam.  The frozen bottleneck codes of both families feed the same four classifiers,
giving Tasks 3 and 4; family-specific denoising variants trained against clean targets under
$20\,\%$ and $40\,\%$ masked-pixel corruption give Task 5; and the maximally activating training
image plus the incoming encoder weight of each bottleneck unit give Task 6.  The principal
empirical finding is that a learned nonlinear code beats the linear PCA code of the same width at
every width measured: the one-hidden code leads PCA by $+0.11$~pp at $k=32$, $+0.05$~pp at $k=64$,
$+0.50$~pp at $k=128$ and $+0.95$~pp at $k=256$.  The strongest single configuration is the
three-hidden encoder at $k=64$ ($98.50\,\%$), which closes $50\,\%$ of the gap between PCA's best
width and Assignment~3's $98.76\,\%$ raw-input baseline.  We further find that the learned code preserves Assignment-3's per-class error profile
rather than introducing new failure modes; that the deepest \emph{classifier} architecture, not the
deepest encoder, is what disables training outright, collapsing to $58.8\text{--}79.5\,\%$ in six
runs; that denoising at $20\,\%$ noise costs no classification accuracy at all --- it gains
$0.16$~pp --- while cutting the error on a corrupted input to $59\,\%$ of the copy-the-input
baseline; and that the inherited
single-epoch stopping rule reports badly underfit models as converged and had to be replaced by a
patience window for classifiers and a relative-plateau window for the autoencoders.
\end{abstract}

\tableofcontents
\newpage

\section{Dataset and Problem Setup}\label{sec:dataset}

\subsection{Data}\label{sec:dataset:data}
The dataset is the five-class MNIST subset used throughout Assignments 3 and 4, held in
\code{Group\_11/}\{train,val,test\} as \code{.jpg} images organised by class.  Each image is
$28\times28$ greyscale normalised to $[0,1]$ by \code{Grayscale()}$\rightarrow$\code{ToTensor()},
then flattened to $\mathbf{x}\in\R^{D}$ with $D=784$ (Eq.~\eqref{eq:flatten}).  Table~\ref{tab:dataset}
gives the split sizes.  The split is exactly balanced, so chance accuracy is exactly $1/C=20\,\%$ on
every split.

\begin{table}[htbp]\centering
\caption{Split sizes and per-class counts of the five-class MNIST subset.  Each class contributes
one fifth of every split, so all splits are balanced; counts are exact, being the split sizes
divided by $C=5$.}\label{tab:dataset}
\begin{tabular}{@{}lrrrrrrr@{}}
\toprule Split & Total & Digit 0 & Digit 4 & Digit 5 & Digit 6 & Digit 7 & Share \\ \midrule
Train & 11\,385 & 2\,277 & 2\,277 & 2\,277 & 2\,277 & 2\,277 & 60.0\,\% \\
Val   & 3\,795  & 759    & 759    & 759    & 759    & 759    & 20.0\,\% \\
Test  & 3\,795  & 759    & 759    & 759    & 759    & 759    & 20.0\,\% \\ \midrule
Total & 18\,975 & 3\,795 & 3\,795 & 3\,795 & 3\,795 & 3\,795 & 100\,\% \\
\bottomrule
\end{tabular}
\end{table}

Tasks 1, 3, 4 and 5 use the training split alone for fitting and for deriving statistics.
Validation accuracy performs \emph{architecture} selection; the bottleneck \emph{width} is selected
by test accuracy, because Assignment~4 asks which width classifies best and reads the answer off
the test scores.  Two consequences are stated up front rather than discovered late.  First, test
accuracy is computed and reported for \emph{every} (architecture, $k$) pair, not once for a single
selected model --- sixteen numbers appear in each of Tables~\ref{tab:task1}, \ref{tab:task3}
and~\ref{tab:task4} --- so the width choice is made over a recorded set rather than over a number
generated for the purpose.  Second, a maximum taken over several test scores is optimistic by
construction; Section~\ref{sec:selbias} measures how optimistic and reports the unbiased
alternative beside it.  This is not the Assignment-3 protocol, so comparisons are of representation
quality rather than of matched training.

\subsection{Notation}\label{sec:dataset:notation}
Table~\ref{tab:notation} fixes the notation for the rest of the report.  Each symbol carries
exactly one meaning throughout: $x$, $z$, $k$, $W$ and $D$ are never redefined, and $D$ is always
the input dimension while $k$ is always the bottleneck width.

\begin{table}[htbp]\centering\small
\caption{Notation.  Each symbol has a single meaning throughout the report.}\label{tab:notation}
\begin{tabular}{@{}p{0.10\linewidth}p{0.82\linewidth}@{}}
\toprule \textbf{Symbol} & \textbf{Meaning} \\ \midrule
$D$ & Input (image) dimension; $D=784=28\times28$. \\
$C$ & Number of classes; $C=5$. \\
$N$ & Number of samples in the split under discussion; $N_{\mathrm{tr}}=11\,385$,
      $N_{\mathrm{val}}=N_{\mathrm{test}}=3\,795$. \\
$\mathcal{S}_{\mathrm{tr}},\mathcal{S}_{\mathrm{val}},\mathcal{S}_{\mathrm{test}}$
    & Training, validation and test sample sets. \\
$y\in\{0,4,5,6,7\}$ & Class label.  Labels are digit values, not indices $0\ldots4$, and are
      printed verbatim in every confusion matrix. \\
$\mathbf{x}$ & One flattened input image, $\mathbf{x}\in\R^{D}$, entries in $[0,1]$. \\
$\mathbf{z}$ & Bottleneck (compressed) representation, $\mathbf{z}\in\R^{k}$, produced by PCA
      projection or by an encoder depending on context. \\
$k$ & Bottleneck width, $k\in\{32,64,128,256\}$. \\
$\mathbf{h}$ & The $400$-wide hidden vector of the three-hidden autoencoder. \\
$\hat{\mathbf{x}}$ & Reconstruction of $\mathbf{x}$, $\hat{\mathbf{x}}\in\R^{D}$; the sigmoid
      output layer keeps it in $[0,1]$ like the target. \\
$\mathbf{W},\mathbf{b}$ & Weight matrix and bias vector.  Superscript $\mathbf{W}^{(l)}$ indexes
      stacked classifier layers; subscripts $e,d$ select an autoencoder's encoder or decoder;
      subscript $e,j$ selects row $j$ of an encoder matrix. \\
$\sigma(\cdot)$ & Logistic sigmoid, Eq.~\eqref{eq:sigmoid}. \\
$\theta$ & All trainable parameters of a model; $\theta_e,\theta_d$ are the encoder and decoder
      subsets and $\theta_{\mathrm{clf}}$ a classifier. \\
$\phi,\psi$ & Encoder and decoder maps; $\hat{\mathbf{x}}=\psi(\phi(\mathbf{x}))$. \\
$\boldsymbol{\mu}$ & Training-set mean image (PCA only). \\
$\bar{\mathbf{x}}$ & Mean-centred image, $\bar{\mathbf{x}}=\mathbf{x}-\boldsymbol{\mu}$. \\
$\boldsymbol{\Sigma}$ & PCA sample covariance matrix, $\boldsymbol{\Sigma}\in\R^{D\times D}$. \\
$\lambda_j,\mathbf{u}_j$ & Eigenvalue and eigenvector of $\boldsymbol{\Sigma}$, ordered
      $\lambda_1\ge\cdots\ge\lambda_D$. \\
$\mathbf{U}_k$ & Matrix of the $k$ leading eigenvectors, $\mathbf{U}_k\in\R^{D\times k}$. \\
$\nu(k)$ & Fraction of total training variance retained by the top-$k$ components. \\
$\rho$ & Fraction of input pixels zeroed by the corruption mask (Task 5 only),
      $\rho\in\{0.20,0.40\}$. \\
$\mathbf{m},\tilde{\mathbf{x}}$ & Binary corruption mask and corrupted input. \\
$\mathbf{s},p_c$ & Logit vector and softmax probability of class $c$. \\
$\mathbf{g}_t$ & Gradient of the objective with respect to $\theta$ at step $t$. \\
$\beta_1,\beta_2,\epsilon_{\mathrm{adam}}$ & Adam's moment decay rates and stabiliser. \\
$\mathcal{L}_{\mathrm{rec}},\mathcal{L}_{\mathrm{CE}},\mathcal{L}_{\mathrm{DAE}}$ &
      Reconstruction, classification and denoising objectives. \\
$E_{\mathrm{rec}}(\mathcal{S})$ & Post-training mean squared reconstruction error on split
      $\mathcal{S}$. \\
$\eta$ & Learning rate; $\eta_{\mathrm{clf}}$ and $\eta_{\mathrm{AE}}$ distinguish the
      classifier and autoencoder rates. \\
$\epsilon_{\mathrm{tol}},\tau_{\mathrm{pat}}$ & Classifier stopping tolerance, $10^{-4}$, and its
      consecutive-epoch patience window, $15$.  Autoencoders instead use a relative plateau rule
      (Section~\ref{sec:arch:opt}) with tolerance $10^{-3}$ and window $50$. \\
$\Acc_{\mathrm{val}},\Acc_{\mathrm{test}}$ & Validation and test accuracy. \\
\bottomrule
\end{tabular}
\end{table}

\subsection{Baseline and metrics}\label{sec:dataset:baseline}
The reference point for every comparison is the best model of Assignment 3, the FCNN
$\mathrm{A\_128\_64\_32}$ ($784\rightarrow128\rightarrow64\rightarrow32\rightarrow5$) trained
with Nesterov-accelerated gradient descent, which reached $\Acc_{\mathrm{val}}=98.84\,\%$ and
$\Acc_{\mathrm{test}}=98.76\,\%$ in 18 epochs.  We quote those numbers from the submitted
Assignment-3 report rather than from the leftover \code{results/} directory of that assignment,
which disagrees with it.  Accuracy, the test confusion matrix and its row-normalised version,
per-class recall, and the mean squared reconstruction error of Eq.~\eqref{eq:recsplit} are the
metrics reported below.  Unless stated otherwise, evaluation is full-batch.

\section{Model Architecture and Implementation}\label{sec:arch}
All models are implemented in PyTorch on a single device selected at startup, CUDA when available
and CPU otherwise, in fp32 throughout.  Mixed precision is off and was not used for any reported
number; the PCA eigendecomposition, the mean subtraction, all reconstruction-error reporting and
the maximally-activating-input search of Task 6 were additionally run in explicit fp32, because an
argmax ordering is not a quantity to expose to reduced-precision reductions.

\subsection{The FCNN classifier}\label{sec:arch:clf}
Every reduced representation in Tasks 1, 3, 4 and 5 is classified by one of the four fully
connected architectures of Table~\ref{tab:archs}.  The set is held \emph{fixed across input
dimensions}: the same shapes are applied to the $784$-dimensional raw input and to the
$k=32,64,128,256$-dimensional reduced representations.  If each $k$ were given architectures
rescaled to its own width, a difference between two widths would be confounded with a difference
in classifier capacity and the experiment would no longer answer which representation classifies
better.  Holding the shapes constant means every width faces identical capacity, so a measured
difference is attributable to the representation.

\begin{table}[htbp]\centering
\caption{The four FCNN classifier architectures, shared by Tasks 1, 3, 4 and 5 at every input
dimension.  Widths funnel monotonically toward the $5$-way output while spanning shallow-wide to
deep-narrow.}\label{tab:archs}
\begin{tabular}{@{}llp{0.40\linewidth}@{}}
\toprule \textbf{Name} & \textbf{Hidden widths} & \textbf{Role} \\ \midrule
3L\_A & $128\to64\to32$ & Shallow, widest capacity. \\
3L\_B & $64\to32\to16$ & Shallow, narrow; capacity-floor probe. \\
4L\_A & $128\to64\to32\to16$ & Mid-depth. \\
5L\_A & $64\to32\to16\to8\to4$ & Deepest, narrowest. \\
\bottomrule
\end{tabular}
\end{table}

A classifier with $L$ hidden layers computes logits $\mathbf{s}\in\R^{C}$ and probabilities $p_c$
by Eqs.~\eqref{eq:clf}--\eqref{eq:ce}.
\begin{align}
\mathbf{s} &= \mathbf{W}^{(L)}\sigma\bigl(\mathbf{W}^{(L-1)}\sigma(\cdots
  \sigma(\mathbf{W}^{(1)}\mathbf{x}+\mathbf{b}^{(1)})\cdots)+\mathbf{b}^{(L-1)}\bigr)
  +\mathbf{b}^{(L)}\in\R^{C}, \label{eq:clf}\\
p_c &= \frac{e^{s_c}}{\sum_{c'=1}^{C}e^{s_{c'}}},\qquad c=1,\ldots,C,\qquad C=5, \label{eq:softmax}\\
\mathcal{L}_{\mathrm{CE}}(\theta_{\mathrm{clf}}) &=
  -\frac{1}{N}\sum_{i=1}^{N}\log p_{y_i}. \label{eq:ce}
\end{align}
Weights are Xavier Glorot uniform with zero biases, drawn under \code{torch.manual\_seed(42)} so
an architecture is reproducible as a unit.  Xavier rather than PyTorch's default Kaiming uniform,
because the funnel widths of Table~\ref{tab:archs} sit in the symmetric regime Xavier was derived
for; the distinction is recorded because Assignment 3's report described Kaiming while its code in
fact used Xavier, a mismatch this report does not repeat.

\subsection{Principal component analysis (Task 1)}\label{sec:arch:pca}
PCA is fitted on $\mathcal{S}_{\mathrm{tr}}$ alone.  Images are flattened, the training mean is
computed once from the training split, every split is centred with that \emph{same} mean, and the
covariance matrix is diagonalised.  The $k$ leading eigenvectors form $\mathbf{U}_k$ and define
the compressed representation of Eq.~\eqref{eq:proj}; Eq.~\eqref{eq:varratio} is the fraction of
total training variance a width-$k$ code retains.  The eigendecomposition is computed in fp32 and
the ranking of eigenvalues is a discrete event, so the top-$k$ choice is deterministic.
\begin{align}
\mathbf{x}_i &= \operatorname{vec}(\mathbf{g}_i)\in\R^{D},\qquad D=28\times28=784,
  \qquad i=1,\ldots,N, \label{eq:flatten}\\
\boldsymbol{\mu} &= \frac{1}{N}\sum_{i=1}^{N}\mathbf{x}_i\in\R^{D},
  \qquad \boldsymbol{\mu}=\boldsymbol{\mu}(\mathcal{S}_{\mathrm{tr}}), \label{eq:mean}\\
\bar{\mathbf{x}}_i &= \mathbf{x}_i-\boldsymbol{\mu}
  \qquad\text{(the training mean, for every split)}, \label{eq:center}\\
\boldsymbol{\Sigma} &= \frac{1}{N}\sum_{i=1}^{N}\bar{\mathbf{x}}_i\bar{\mathbf{x}}_i^{\top}
  \in\R^{D\times D},\qquad \boldsymbol{\Sigma}=\boldsymbol{\Sigma}^{\top}\succeq0, \label{eq:cov}\\
\boldsymbol{\Sigma}\mathbf{u}_j &= \lambda_j\mathbf{u}_j,\qquad
  \lambda_1\ge\cdots\ge\lambda_D\ge0,\qquad
  \langle\mathbf{u}_j,\mathbf{u}_{j'}\rangle=\delta_{jj'}, \label{eq:eig}\\
\mathbf{U}_k &= [\mathbf{u}_1\;\mathbf{u}_2\;\cdots\;\mathbf{u}_k]\in\R^{D\times k},
  \qquad \mathbf{z}_i=\mathbf{U}_k^{\top}\bar{\mathbf{x}}_i\in\R^{k}, \label{eq:proj}\\
\nu(k) &= \frac{\sum_{j=1}^{k}\lambda_j}{\sum_{j=1}^{D}\lambda_j}\in[0,1]
  \qquad\text{(non-decreasing in $k$)}. \label{eq:varratio}
\end{align}
where $\mathbf{g}_i\in[0,1]^{28\times28}$ is the greyscale image of Eq.~\eqref{eq:flatten}.

\subsection{Autoencoders (Tasks 2, 3 and 4)}\label{sec:arch:ae}
Two encoder/decoder families are trained.  In both, the bottleneck layer is \emph{linear} by
mandate and every other hidden layer as well as the output layer uses the logistic sigmoid of
Eq.~\eqref{eq:sigmoid}.  Logistic is used throughout rather than $\tanh$, because $\tanh$ forces
targets into $[-1,1]$, putting reconstruction error outside pixel units and making the Task-2 and
Task-5 numbers incomparable with each other.
\begin{equation}\label{eq:sigmoid}
\sigma(a)=\frac{1}{1+e^{-a}}\in(0,1),\qquad \sigma'(a)=\sigma(a)\bigl(1-\sigma(a)\bigr).
\end{equation}

\paragraph{One-hidden family ($784\rightarrow k\rightarrow784$).}
The bottleneck is the only hidden layer, so it is linear, and the decoder is a single sigmoid
output layer, giving Eq.~\eqref{eq:ae1full}.  The parameter count is $2Dk+D+k$.
\begin{align}
\mathbf{z} &= \phi(\mathbf{x})=\mathbf{W}_e\mathbf{x}+\mathbf{b}_e\in\R^{k},
  \qquad \theta_e=\{\mathbf{W}_e\in\R^{k\times D},\mathbf{b}_e\}, \label{eq:ae1enc}\\
\hat{\mathbf{x}} &= (\psi\circ\phi)(\mathbf{x})=\sigma\bigl(\mathbf{W}_d
  (\mathbf{W}_e\mathbf{x}+\mathbf{b}_e)+\mathbf{b}_d\bigr)\in\R^{D}, \label{eq:ae1full}\\
\psi(\mathbf{z}) &= \sigma(\mathbf{W}_d\mathbf{z}+\mathbf{b}_d),\qquad
  \theta_d=\{\mathbf{W}_d\in\R^{D\times k},\mathbf{b}_d\}. \label{eq:ae1dec}
\end{align}

\paragraph{Three-hidden family ($784\rightarrow400\rightarrow k\rightarrow400\rightarrow784$).}
Two sigmoid layers of width $400$ bracket the same linear bottleneck.  The $400$-wide layers are
the assignment's own prescription, not a tuned width.  The parameter count is
$628\,784+801k$, dominated by the two $784\leftrightarrow400$ matrices.
\begin{align}
\mathbf{h} &= \sigma\bigl(\mathbf{W}_e^{(1)}\mathbf{x}+\mathbf{b}_e^{(1)}\bigr)\in\R^{400},
  \qquad \mathbf{z}=\mathbf{W}_e^{(2)}\mathbf{h}+\mathbf{b}_e^{(2)}\in\R^{k},
  \label{eq:ae3enc}\\
\mathbf{h}' &= \sigma\bigl(\mathbf{W}_d^{(1)}\mathbf{z}+\mathbf{b}_d^{(1)}\bigr)\in\R^{400},
  \qquad \hat{\mathbf{x}}=\sigma\bigl(\mathbf{W}_d^{(2)}\mathbf{h}'+\mathbf{b}_d^{(2)}\bigr)
  \in\R^{D}. \label{eq:ae3dec}
\end{align}

\begin{table}[htbp]\centering
\caption{Exact trainable-parameter counts of the two autoencoder families.  The three-hidden
family carries a large fixed cost from the two $784\leftrightarrow400$ matrices, independent of
$k$.}\label{tab:aeparams}
\begin{tabular}{@{}crrrr@{}}
\toprule Family & $k=32$ & $k=64$ & $k=128$ & $k=256$ \\ \midrule
One-hidden   & 50\,992 & 101\,200 & 201\,616 & 402\,448 \\
Three-hidden & 654\,416 & 680\,048 & 731\,312 & 833\,840 \\
\bottomrule
\end{tabular}
\end{table}

Both families minimise the mean squared reconstruction error of Eq.~\eqref{eq:mse} with Adam.  Two
points of definition matter.  First, $E_{\mathrm{rec}}$ is evaluated \emph{after} training, as a
full-batch pass over the split, and is not accumulated during training, so the numbers in
Table~\ref{tab:recon} are the error of the final weights and of nothing else.  Second,
$\hat{\mathbf{x}}$ comes from a sigmoid and therefore lies in $[0,1]$, exactly the range of the
targets, so the error is in pixel-intensity-squared units and needs no rescaling.
\begin{align}
\mathcal{L}_{\mathrm{rec}}(\theta) &= \frac{1}{N\,D}\sum_{i=1}^{N}\sum_{j=1}^{D}
  \bigl(x_{ij}-\hat{x}_{ij}\bigr)^{2},\qquad \theta=\theta_e\cup\theta_d, \label{eq:mse}\\
E_{\mathrm{rec}}(\mathcal{S}) &= \frac{1}{|\mathcal{S}|\,D}
  \sum_{\mathbf{x}\in\mathcal{S}}\bigl\|\mathbf{x}-\hat{\mathbf{x}}(\mathbf{x})\bigr\|_2^{2},
  \qquad \mathcal{S}\in\{\mathcal{S}_{\mathrm{tr}},\mathcal{S}_{\mathrm{val}},
  \mathcal{S}_{\mathrm{test}}\}. \label{eq:recsplit}
\end{align}

\subsection{Denoising autoencoders (Task 5)}\label{sec:arch:dae}
The denoising autoencoders are structurally identical to the one-hidden autoencoder of
Eq.~\eqref{eq:ae1full}; what changes is the training distribution, not the architecture.  Each
input is corrupted by a fresh per-sample binary mask and the reconstruction is scored against the
\emph{clean} image.  Masking corruption is used rather than additive Gaussian noise: it models
salt-like pixel dropout, keeps the corrupted input inside $[0,1]$ so the target distribution is
unchanged, and keeps the sigmoid output well matched to that distribution.  Because the target is
clean and the input is not, the bottleneck cannot carry the corruption forward and is obliged to
discard it, which is the mechanism by which a denoising code differs from a plain code of the
same width.
\begin{align}
\tilde{\mathbf{x}} &= \mathbf{x}\odot\mathbf{m},\qquad
  m_j\sim\mathrm{Bernoulli}(1-\rho)\ \text{i.i.d.},\qquad \rho\in\{0.20,\,0.40\}, \label{eq:noise}\\
\mathcal{L}_{\mathrm{DAE}}(\theta) &= \frac{1}{N\,D}\sum_{i=1}^{N}\sum_{j=1}^{D}
  \Bigl(x_{ij}-\bigl[\psi(\phi(\tilde{\mathbf{x}}_i))\bigr]_{j}\Bigr)^{2}. \label{eq:dae}
\end{align}

\subsection{Weight visualisation (Task 6)}\label{sec:arch:wv}
For each unit $j$ of a one-hidden encoder we visualise two objects: the training image that
maximally activates that unit, and the vector of weights feeding it, as in Eq.~\eqref{eq:maxact}.
For the one-hidden family the encoder is a single \code{nn.Linear(784,$k$)}, so its row $j$
\emph{is} the $784\rightarrow j$ connection: $[\mathbf{W}_e]_{j,:}$ is a $784$-dimensional vector
reshapeable to $28\times28$ and displayed with no further processing.  The weight vector is what
the unit \emph{looks for}; the maximally activating image is the training sample on which that
preference is most strongly expressed.
\begin{equation}\label{eq:maxact}
\mathbf{x}^{*}_{j}=\arg\max_{\mathbf{x}\in\mathcal{S}_{\mathrm{tr}}}[\phi(\mathbf{x})]_{j},
\qquad \mathbf{w}_{e,j}^{\top}=[\mathbf{W}_e]_{j,:}\in\R^{D}.
\end{equation}

\subsection{Optimisation, stopping and reproducibility}\label{sec:arch:opt}
All training uses Adam with $\beta_1=0.9$, $\beta_2=0.999$ and $\epsilon_{\mathrm{adam}}=10^{-8}$.
Assignment 4 mandates Adam for the autoencoders; the same optimiser is used for every classifier,
so the four tasks differ in representation only, never in optimiser.  Gradients are taken over the
full training batch throughout.  This differs from Assignment 3's best model, which was trained at
batch size $1$; the consequence is recorded in Section~\ref{sec:selbias}, and it is one reason the
two assignments' accuracies should not be read as a controlled comparison.
\begin{align}
\bar{\mathbf{g}}_t &= \beta_1\bar{\mathbf{g}}_{t-1}+(1-\beta_1)\mathbf{g}_t,
  \qquad \overline{\mathbf{g}^2}_t=\beta_2\overline{\mathbf{g}^2}_{t-1}+(1-\beta_2)\mathbf{g}_t^{2},
  \label{eq:adam}\\[2pt]
\hat{\mathbf{g}}_t &= \frac{\bar{\mathbf{g}}_t}{1-\beta_1^{t}},\qquad
  \widehat{\mathbf{g}^2}_t=\frac{\overline{\mathbf{g}^2}_t}{1-\beta_2^{t}},\qquad
  \theta\leftarrow\theta-\eta\frac{\hat{\mathbf{g}}_t}{\sqrt{\widehat{\mathbf{g}^2}_t}
  +\epsilon_{\mathrm{adam}}},
\end{align}
where $\mathbf{g}_t=\nabla_{\theta}\mathcal{L}$ is the gradient of the objective of the run at
step $t$.  Learning rates follow Assignment 3 where a value already exists
($\eta_{\mathrm{clf}}=0.01$, $\eta_{\mathrm{AE}}=0.001$); both were validated by a pilot on this
dataset, because full-batch descent at $\eta=0.001$ was measured to sit at exactly $20\,\%$
validation accuracy --- chance --- which is the same full-batch failure Assignment 3 recorded and
which this assignment must not reproduce.

\paragraph{Stopping criterion.}
The inherited rule --- stop when the single-epoch loss change falls below
$\epsilon_{\mathrm{tol}}=10^{-4}$ --- was re-examined, found unsound, and replaced.  A
single-epoch difference cannot distinguish genuine settlement from an optimiser oscillating inside
a plateau: on this dataset the 5L\_A classifier stalls with loss $1.6089$ at epoch 13 and then
\emph{rises} monotonically ($1.6091$, $1.6096$, $1.6102$, $1.6108$, $1.6111$), yet Adam's momentum
transiently produces a one-epoch $|\Delta L|$ of order $10^{-5}$, under the tolerance, so the naive
rule declares convergence while validation accuracy is exactly $20.00\,\%$ --- chance.  We
therefore stop only when the criterion holds for $\tau_{\mathrm{pat}}=15$ \emph{consecutive}
epochs; a genuine minimum sustains the small difference, an oscillation does not.  A budget of
$10\,000$ epochs bounds every run.

The classifiers keep that rule verbatim.  The autoencoders do not, and cannot: a mean squared
reconstruction error at a scale of $10^{-3}$ and a cross-entropy at a scale of $1.6$ cannot share
one absolute tolerance, since the same $10^{-4}$ delta is a large relative step for the former and
a negligible one for the latter.  Both autoencoder families are therefore stopped by a
\emph{relative plateau} rule instead: training halts once the monotone best loss has failed to
improve by more than $10^{-3}$ in relative terms over a window of $50$ epochs.  The window acts as
the patience of the rule and is immune to the epoch-to-epoch jitter of the masked corruption, which
is what makes it usable for the denoising runs of Task~5.  Table~\ref{tab:stopconv} records the rule
and the stopping epoch of every run, so no stopping behaviour in this report rests on prose.

\begin{table}[htbp]\centering\small
\caption{Stopping rule and epoch at which each autoencoder stopped.  All ten autoencoder and
denoising runs are stopped by the relative-plateau rule and all stop before the $10\,000$-epoch
budget; none runs to the cap.  The classifiers, by contrast, are stopped by the absolute
$10^{-4}$ rule after a mean of $293$ epochs (range $144$--$731$) across the $48$ runs that are not
degenerate.}\label{tab:stopconv}
\begin{tabular}{@{}llrr@{}}
\toprule Run & Rule & Stopped at epoch & Budget \\ \midrule
1-hidden, $k=32$   & relative plateau & 3\,470 & 10\,000 \\
1-hidden, $k=64$   & relative plateau & 4\,751 & 10\,000 \\
1-hidden, $k=128$  & relative plateau & 5\,594 & 10\,000 \\
1-hidden, $k=256$  & relative plateau & 6\,555 & 10\,000 \\
3-hidden, $k=32$   & relative plateau & 7\,198 & 10\,000 \\
3-hidden, $k=64$   & relative plateau & 7\,763 & 10\,000 \\
3-hidden, $k=128$  & relative plateau & 7\,835 & 10\,000 \\
3-hidden, $k=256$  & relative plateau & 7\,413 & 10\,000 \\ \midrule
Denoising, $\rho=0.20$ & relative plateau & 2\,130 & 10\,000 \\
Denoising, $\rho=0.40$ & relative plateau & 1\,946 & 10\,000 \\ \bottomrule
\end{tabular}
\end{table}

Degeneracy is judged by an absolute floor rather than by chance.  A run below $85\,\%$ validation
accuracy is flagged \emph{degenerate} and excluded from selection; six runs trip it, all of the
5L\_A architecture.  The floor sits between the two clusters actually observed --- genuine results
at $97\text{--}99\,\%$, 5L\_A failures at $58.8$, $59.6$, $77.9$ and $78.3\,\%$ --- which a fixed
multiple of chance cannot separate, since $77.9\,\%$ is already nearly four times chance.  It is a
dataset-specific threshold, and is stated as such.

\paragraph{Reproducibility.}
\code{torch.manual\_seed(42)} governs every weight initialisation and every shuffle, so runs are
reproducible in initialisation and in data order.  They are \emph{not} bitwise reproducible: GPU
reductions have a non-deterministic accumulation order, and making them deterministic would need
\code{torch.use\_deterministic\_algorithms(True)} with \code{cudnn.deterministic=True}, which was
not enabled.  We claim the weaker statement, that a rerun lands on the same model and the same
trajectory up to floating-point reordering.

\paragraph{Checkpointing and progress.}
Every run checkpoints model state, optimiser state, epoch counter and full metric history to
\code{<run>.tmp} and atomically renames it onto \code{<run>.pt} every 10 epochs, plus an
unconditional final checkpoint written \emph{after} the reported metrics are computed, so a crash
costs at most ten epochs, a truncated file can never be mistaken for a loadable one, and the
weights on disk are exactly the weights the report describes.  Resume restores model, optimiser
\emph{and} the full convergence-watch state --- the mode, the window, the minimum-epoch setting and
the windowed best-loss history, not merely the patience streak.  Restoring less would be a silent
correctness bug: a resumed autoencoder would fall back to the constructor's default mode and
re-evaluate the plateau test from an empty history, and could then stop far earlier than the rule
intends.  Progress lines and status snapshots use a rolling 50-epoch rate window
rather than a whole-run mean, because sustained load on a $4\,$GB laptop GPU thermally throttles
and a mean-rate ETA only ever becomes more optimistic as the run proceeds.

\begin{lstlisting}[caption={Autoencoder construction (models.py).  The bottleneck is linear in both
families; the three-hidden family brackets it with 400-wide sigmoid layers as the assignment
prescribes.  The same single \code{nn.Linear(784,$k$)} encoder is reused, with a masking
corruption of Eq.~\eqref{eq:noise} applied to its input, as the denoising autoencoder of Task 5.},
label={lst:ae}]
if kind == "1hidden":                       # 784 -> k -> 784
    encoder = nn.Sequential(nn.Linear(input_dim, bottleneck))
    decoder = nn.Sequential(nn.Linear(bottleneck, input_dim), nn.Sigmoid())
elif kind == "3hidden":                     # 784 -> 400 -> k -> 400 -> 784
    encoder = nn.Sequential(nn.Linear(input_dim, 400), nn.Sigmoid(),
                            nn.Linear(400, bottleneck))   # linear bottleneck
    decoder = nn.Sequential(nn.Linear(bottleneck, 400), nn.Sigmoid(),
                            nn.Linear(400, input_dim), nn.Sigmoid())
\end{lstlisting}

\section{Task 1: Dimensionality Reduction with PCA}\label{sec:task1}
For each of $k=32,64,128,256$ the training split is projected onto the top-$k$ eigenvectors of
$\boldsymbol{\Sigma}$ per Eqs.~\eqref{eq:center}--\eqref{eq:proj}, and all four classifiers of
Table~\ref{tab:archs} are trained on the resulting code.  Validation accuracy over every
(architecture, $k$) pair selects one architecture per $k$; the test accuracy and confusion matrix
are then computed for every architecture at every width, not only for the selected one, because
Assignment~4 selects the width on test accuracy
(Section~\ref{sec:dataset:data}).  Table~\ref{tab:task1} gives the validation surface and
Figure~\ref{fig:task1:heat} the corresponding $16$-cell test grid.

\begin{table}[htbp]\centering
\caption{Task 1: validation accuracy (\%) of each classifier on the PCA code at each width.  Rows
are the four architectures of Table~\ref{tab:archs}, columns are $k$.  The best entry of each
column selects the architecture evaluated on the test split.}\label{tab:task1}
\begin{tabular}{@{}lrrrr@{}}
\toprule Architecture & $k=32$ & $k=64$ & $k=128$ & $k=256$ \\ \midrule
3L\_A & 98.39 & 98.37 & 98.10 & 97.05 \\
3L\_B & 98.55 & 98.21 & 98.00 & 97.29 \\
4L\_A & 98.50 & 98.02 & 97.89 & 96.89 \\
5L\_A & 97.94 & 79.45 & 77.87 & 96.86 \\
\bottomrule
\end{tabular}
\end{table}

\paragraph{Eigen-spectrum and variance retained.}
Figure~\ref{fig:task1:spectrum} shows the eigenvalue spectrum of $\boldsymbol{\Sigma}$ and
Figure~\ref{fig:task1:variance} the retained variance $\nu(k)$ against $k$.  The two are the same
curve at different resolutions and together they decide how much of the $784$-dimensional image a
width-$k$ code can carry at all.  The shape --- a few large eigenvalues followed by a long shallow
tail --- is the standard MNIST result: the leading components span the average stroke direction and
the global intensity variation, which is why $\nu(k)$ rises steeply at small $k$ and then flattens.
The flattening is the part that matters, because it bounds what PCA can deliver: beyond roughly
$k=64$ each additional component buys a small and decreasing share of variance, so the code stops
getting richer much faster than it stops getting wider.

\begin{figure}[htbp]\centering
\maybeinc[width=0.84\linewidth]{task1/07_representations/task1_eigen_spectrum.png}
\caption{Task 1: eigenvalue spectrum of the training covariance matrix $\boldsymbol{\Sigma}$,
logarithmic scale.  The steep drop-off after the first few components is why $\nu(k)$ rises quickly
and then flattens.}\label{fig:task1:spectrum}
\end{figure}

\begin{figure}[htbp]\centering
\maybeinc[width=0.84\linewidth]{task1/07_representations/task1_variance_retained.png}
\caption{Task 1: fraction $\nu(k)$ of total training variance retained by the top-$k$ principal
components (Eq.~\eqref{eq:varratio}), annotated at each required width.  The curve saturates, so
the marginal information gained per additional component falls as $k$ grows; this is the linear
analogue of the diminishing returns the autoencoder codes show in Section~\ref{sec:task2}.}
\label{fig:task1:variance}
\end{figure}

\paragraph{Accuracy versus width.}
Figure~\ref{fig:task1:acc} traces validation accuracy against $k$ for the selected architecture at
each width.  The sign of the trend is unambiguous: it falls, monotonically in the test column and
almost so in the validation column.  That is worth a paragraph of its own rather than a reading
guide, because a falling curve against a rising variance curve is the sort of figure that is
misread.

\begin{figure}[htbp]\centering
\maybeinc[width=0.84\linewidth]{task1/08_comparisons/task1_accuracy_vs_dimension.png}
\caption{Task 1: validation accuracy of the selected classifier against the PCA code width $k$.
Compared directly with Figure~\ref{fig:task3:acc} in Section~\ref{sec:results}.}
\label{fig:task1:acc}
\end{figure}

\paragraph{Best configuration.}
The best PCA width is $k=32$, selected by test accuracy as Assignment~4 prescribes, with the
3L\_B classifier at $\Acc_{\mathrm{val}}=98.55\,\%$ and
$\Acc_{\mathrm{test}}=98.23\,\%$ ($3728/3795$).  The confusion structure in
Figure~\ref{fig:task1:cm} is inherited rather than new: the dominant off-diagonal mass lies on
the $5\leftrightarrow6$ pair, then $4\leftrightarrow7$ and $5\leftrightarrow7$, the same ordering
Assignment~3 reported on the raw input.  That is the first indication the linear code preserves the
structure of the original errors rather than introducing a different failure mode.

\paragraph{Why accuracy falls as width rises.}
Test accuracy falls from $98.2345\,\%$ at $k=32$ to $97.2332\,\%$ at $k=256$ while retained variance
\emph{rises} from $0.759$ to $0.979$.  The representations are not deteriorating; the classifiers
are.  The four architectures are deliberately held fixed across widths
(Section~\ref{sec:arch:clf}), so as $k$ grows the first layer's input dimension grows from $32$ to
$256$ while its width stays fixed and it becomes progressively more ill-conditioned.  A wider code
is being fed to a badly shaped first layer, not to a bigger one.  The reading is therefore that PCA
already carries the class information by about $k=32$ and that everything above it is variance the
fixed-width input layer cannot exploit.  This is a finding about the architecture policy rather
than about PCA, and it has a consequence worth flagging for anyone reusing these results: the
\emph{falling} part of the curve is a capacity artefact of holding the widths fixed, and the
\emph{level} at $k=32$ is the quantity comparable to the autoencoder codes.  It also means the
width axis of every figure in this report carries two variables at once --- representation width
and first-layer conditioning --- so a difference along it is never attributable to width alone.

\begin{figure}[htbp]\centering
\maybeinc[width=0.78\linewidth]{task1/03_confusion_matrices/task1_k32_3L_B_confusion_test.png}
\caption{Task 1: test confusion matrix of the selected configuration, $k=32$ with the 3L\_B
classifier.  Rows are true digits, columns predicted digits, and both are digit values rather than
class indices.  The $5\leftrightarrow6$ pair carries the largest off-diagonal mass.}
\label{fig:task1:cm}
\end{figure}

\section{Task 2: Image Reconstruction with Autoencoders}\label{sec:task2}
Eight autoencoders are trained --- two families $\times$ four bottleneck widths --- each minimising
Eq.~\eqref{eq:mse} with Adam at $\eta_{\mathrm{AE}}=0.001$.  Table~\ref{tab:recon} reports
$E_{\mathrm{rec}}$ on all three splits after training and Figure~\ref{fig:task2:err} plots the same
quantity against $k$.  Two gaps are informative in their own right: the vertical gap between
$E_{\mathrm{rec}}(\mathcal{S}_{\mathrm{tr}})$ and $E_{\mathrm{rec}}(\mathcal{S}_{\mathrm{test}})$,
the generalisation gap of the reconstruction, and the horizontal gap between families at the same
$k$, the price of the extra depth.

\begin{table}[htbp]\centering
\caption{Task 2: post-training mean squared reconstruction error $E_{\mathrm{rec}}$
(Eq.~\eqref{eq:recsplit}) for the eight autoencoders.  The last column is the reconstruction
generalisation gap $E_{\mathrm{rec}}(\mathcal{S}_{\mathrm{test}})-E_{\mathrm{rec}}
(\mathcal{S}_{\mathrm{tr}})$.}\label{tab:recon}
\begin{tabular}{@{}llrrrr@{}}
\toprule Family & $k$ & Train & Val & Test & Gap \\ \midrule
One-hidden    & 32   & 0.008544 & 0.009175 & 0.009346 & 0.000802 \\
One-hidden    & 64   & 0.003137 & 0.003705 & 0.003748 & 0.000611 \\
One-hidden    & 128  & 0.001299 & 0.001684 & 0.001707 & 0.000408 \\
One-hidden    & 256  & 0.000641 & 0.000909 & 0.000926 & 0.000285 \\
Three-hidden  & 32   & 0.003209 & 0.006427 & 0.006537 & 0.003328 \\
Three-hidden  & 64   & 0.001948 & 0.003795 & 0.003859 & 0.001911 \\
Three-hidden  & 128  & 0.001478 & 0.002737 & 0.002772 & 0.001294 \\
Three-hidden  & 256  & 0.001330 & 0.002317 & 0.002346 & 0.001016 \\
\bottomrule
\end{tabular}
\end{table}

\begin{figure}[htbp]\centering
\maybeinc[width=0.84\linewidth]{task2/05_reconstruction_error/task2_recon_error_all.png}
\caption{Task 2: post-training reconstruction error $E_{\mathrm{rec}}$ against bottleneck width $k$
for both families and all three splits.  The $y$-axis is shared, so the separation between families
is directly comparable.}\label{fig:task2:err}
\end{figure}

\paragraph{Reading the reconstruction error.}
The first check in Figure~\ref{fig:task2:err} is whether the three curves of a family move together
or fan out.  If train, val and test coincide at small $k$ and separate at large $k$, the autoencoder
is interpolating the training set rather than learning a code, and the point of separation is the
width at which the code stops generalising; if they stay together throughout, the bottleneck is
acting as a sufficient statistic even at $k=256$ and the family is not overfitting.  The second
check is the family separation: because the three-hidden family has a fixed $\approx628\,784$
parameters regardless of $k$ (Table~\ref{tab:aeparams}), it can afford an extremely nonlinear map
for a narrow code, and the question is whether that nonlinearity buys a lower $E_{\mathrm{rec}}$ at
equal $k$ or merely a lower training error.  A family gap that opens on the training curve but
closes on the test curve is memorisation, and is reported as such.

\paragraph{Reconstructed images.}
Figure~\ref{fig:task2:recon1} shows, for the test split, one image per class alongside its
reconstruction; these are the figures that convert Table~\ref{tab:recon} into something readable,
because a test error on its own says nothing about whether a $4$ was recovered as a $4$.  What to
look for, and what Section~\ref{sec:obs} reports, is the \emph{order in which} digits survive the
bottleneck: digits with large closed strokes and high stroke coverage generally survive at $k=32$,
while digits made of thin separated curves generally do not, because the coded features
distinguishing them are exactly the ones a narrow linear bottleneck is least able to preserve.  A
third item is the character of the failure at $k=32$: a blurred, class-mean-like reconstruction is a
capacity failure, whereas one that is sharp but is of the wrong digit is a code-allocation failure,
and the two carry different implications for Task 3.

\begin{figure}[htbp]\centering
\maybeinc[width=0.95\linewidth]{task2/04_reconstructions/task2_1hidden_k32_recon_test.png}
\caption{Task 2, one-hidden family, $k=32$: originals (top) and reconstructions (bottom) for one
image per class from the test split.  Read together with the $k=64,128,256$ panels of the same
folder and with the corresponding three-hidden panels archived in the
appendix.}\label{fig:task2:recon1}
\end{figure}

\section{Task 3: Classification from the One-Hidden Autoencoder}\label{sec:task3}
For each of the four one-hidden autoencoders the encoder $\phi$ is frozen and applied to all three
splits, producing codes $\mathbf{z}\in\R^{k}$ per Eq.~\eqref{eq:ae1enc}.  The four classifiers of
Table~\ref{tab:archs} are then trained on each code at input dimension $k$, exactly as in Task 1.
Because the autoencoders were already trained in Task 2, no retraining occurs here; this matters
for reproducibility, since a retrained autoencoder would give a different code and therefore
different accuracy for reasons unrelated to the quality of the code.

\begin{table}[htbp]\centering
\caption{Task 3: validation accuracy (\%) of each classifier on the frozen one-hidden code at each
width.  The best entry of each column selects the architecture evaluated on the test
split.}\label{tab:task3}
\begin{tabular}{@{}lrrrr@{}}
\toprule Architecture & $k=32$ & $k=64$ & $k=128$ & $k=256$ \\ \midrule
3L\_A & 98.29 & 98.34 & 98.18 & 98.23 \\
3L\_B & 98.29 & 98.26 & 98.21 & 97.89 \\
4L\_A & 98.08 & 98.21 & 97.81 & 97.89 \\
5L\_A & 97.73 & 77.87 & 59.63 & 58.81 \\
\bottomrule
\end{tabular}
\end{table}

\begin{figure}[htbp]\centering
\maybeinc[width=0.84\linewidth]{task3/08_comparisons/task3_accuracy_vs_bottleneck.png}
\caption{Task 3: validation accuracy against bottleneck width $k$ for the frozen one-hidden code.
The comparison that matters is against Figure~\ref{fig:task1:acc}: same four classifiers, same
widths, different code.}\label{fig:task3:acc}
\end{figure}

The best one-hidden configuration is $k=32$ with the 3L\_A classifier, at
$\Acc_{\mathrm{val}}=98.29\,\%$ and $\Acc_{\mathrm{test}}=98.34\,\%$ ($3732/3795$).  Its test
confusion matrix is Figure~\ref{fig:task3:cm}.  Both claims hold.  The $5\leftrightarrow6$
confusion remains the dominant off-diagonal mass, and the per-class recall ordering is unchanged
from PCA and from the raw input, so the learned code inherited the error structure of the problem
rather than reshaping it.  It should be recorded, however, that $k=32$ --- not $k=256$ --- is the
best one-hidden width: at $k=256$ the one-hidden code reaches $98.18\,\%$, which is below PCA's
best width ($98.23\,\%$) rather than above it.

\begin{figure}[htbp]\centering
\maybeinc[width=0.78\linewidth]{task3/03_confusion_matrices/task3_k32_3L_A_confusion_test.png}
\caption{Task 3: test confusion matrix of the selected configuration, $k=32$ with the 3L\_A
classifier.  Rows are true digits, columns predicted digits.  Compare with
Figure~\ref{fig:task1:cm} at identical labels.}
\label{fig:task3:cm}
\end{figure}

\section{Task 4: Classification from the Three-Hidden Autoencoder}\label{sec:task4}
Task 4 repeats Task 3 with the frozen encoder of the three-hidden family,
$\mathbf{z}=\mathbf{W}_e^{(2)}\sigma(\mathbf{W}_e^{(1)}\mathbf{x}+\mathbf{b}_e^{(1)})+
\mathbf{b}_e^{(2)}$.  The assignment statement calls this the ``2-hidden'' autoencoder, counting
only the layers beyond the bottleneck that carry a nonlinearity; we call it three-hidden throughout
because the object actually built has three hidden layers of widths $400$, $k$ and $400$, while the
one-hidden object of Task 3 has exactly one, of width $k$.  Both names refer to the same two models,
and the naming is stated so that the count never drifts between the code, this report and the
figures.

\begin{table}[htbp]\centering
\caption{Task 4: validation accuracy (\%) of each classifier on the frozen three-hidden code at
each width.  Degenerate runs are reported verbatim rather than blanked, and are excluded from
selection by the absolute floor of Section~\ref{sec:arch:opt}.}\label{tab:task4}
\begin{tabular}{@{}lrrrr@{}}
\toprule Architecture & $k=32$ & $k=64$ & $k=128$ & $k=256$ \\ \midrule
3L\_A & 98.55 & 98.37 & 98.31 & 98.13 \\
3L\_B & 98.21 & 98.60 & 98.16 & 97.89 \\
4L\_A & 98.39 & 98.60 & 98.05 & 98.34 \\
5L\_A & 98.00 & 97.92 & 98.08 & 78.34 \\
\bottomrule
\end{tabular}
\end{table}


\paragraph{Depth does help here, against the initial expectation.}
Figure~\ref{fig:task4:acc} is the direct counterpart of Figure~\ref{fig:task3:acc}: both families
emit a $\R^{k}$ bottleneck at the same $k$, the classifiers and the protocol are identical, and the
difference is the two $400$-wide sigmoid layers in this encoder.  The expected reading was that
accuracy at matched $k$ would be no better than Task~3's and plausibly worse.  It is not what
happened: the three-hidden code wins at three of the four widths, by $+0.08$ pp at $k=32$,
$+0.42$ pp at $k=64$ and $+0.08$ pp at $k=128$, and loses only at $k=256$, by $-0.37$ pp.  The
wider widths are where it degrades, not the narrower ones, so the extra depth helps until the
encoder is asked to pass a wide code through too much nonlinearity.  Its best result,
$98.50\,\%$ at $k=64$, is the highest of any configuration in this study.

The collapse visible in the bottom row of Table~\ref{tab:task4} is a separate phenomenon and must
not be attributed to the encoder.  It is the single degenerate run of that task --- 5L\_A at
$k=256$, $78.34\,\%$ --- and the same 5L\_A classifier is healthy at every other width
($98.00$, $97.92$ and $98.08\,\%$) on that same encoder.  More strikingly, at $k=64$ and $k=128$ the
three-hidden encoder beats the one-hidden encoder at \emph{every} classifier depth, including
5L\_A: at $k=64$, $97.92\,\%$ against $77.87\,\%$, and at $k=128$, $98.08\,\%$ against
$59.63\,\%$.  The one place the one-hidden encoder wins at these widths is 3L\_B at $k=128$
($98.21$ against $98.16\,\%$), a $0.05$~pp difference.  Whatever is breaking the deepest classifier
is therefore not made worse by the deeper encoder, and is not cured by it either.

\begin{figure}[htbp]\centering
\maybeinc[width=0.84\linewidth]{task4/08_comparisons/task4_accuracy_vs_bottleneck.png}
\caption{Task 4: validation accuracy against bottleneck width $k$ for the frozen three-hidden code.
Read against Figure~\ref{fig:task3:acc} at matched $k$.}\label{fig:task4:acc}
\end{figure}

The best three-hidden configuration is $k=64$ with the 4L\_A classifier, at
$\Acc_{\mathrm{val}}=98.60\,\%$ and $\Acc_{\mathrm{test}}=98.50\,\%$ ($3738/3795$); its test
confusion matrix is Figure~\ref{fig:task4:cm}, and it shows the same dominant
$5\leftrightarrow6$ confusion as Figures~\ref{fig:task1:cm} and~\ref{fig:task3:cm}.  This is the
best configuration measured anywhere in the assignment.  Measured against Assignment~3, it closes
half of the gap that PCA's best width leaves open: PCA sits $0.5255$~pp below the $98.76\,\%$
baseline and this configuration sits $0.2620$~pp below it.

\begin{figure}[htbp]\centering
\maybeinc[width=0.78\linewidth]{task4/03_confusion_matrices/task4_k64_4L_A_confusion_test.png}
\caption{Task 4: test confusion matrix of the selected configuration, $k=64$ with the 4L\_A
classifier.  Compare with Figures~\ref{fig:task1:cm} and~\ref{fig:task3:cm}.}
\label{fig:task4:cm}
\end{figure}

\section{Task 5: Denoising Autoencoders}\label{sec:task5}
The one-hidden autoencoder is retrained twice with the corruption of Eq.~\eqref{eq:noise} at
$\rho=0.20$ and $\rho=0.40$, targeting the clean image per Eq.~\eqref{eq:dae}.  The bottleneck
width is the one selected in Task 3, $k=32$, because the assignment fixes it that way and
because it is the only width at which the denoising code is compared against a known plain-code
baseline rather than against itself.  Table~\ref{tab:task5clf} reports the classification
accuracy of both codes, alongside the reconstruction errors.

\begin{table}[htbp]\centering
\caption{Task 5: classification accuracy from the frozen denoising bottleneck
representation against the plain one-hidden result at the same $k$ (from
Table~\ref{tab:task3}), together with the post-training reconstruction error of
each model on the test split and the Assignment-3 baseline.}
\label{tab:task5clf}
\begin{tabular}{@{}lrrrr@{}}
\toprule Model & Val acc. (\%) & Test acc. (\%) & Test $E_{\mathrm{rec}}$ & $\Delta$ vs. plain \\
\midrule Plain one-hidden, $k=32$, 3L\_A & 98.29 & 98.34 & 0.009346 & --- \\
Denoising, $\rho=0.20$, 3L\_A & 98.50 & 98.50 & 0.015053 & $+0.16$ \\
Denoising, $\rho=0.40$, 3L\_A & 98.60 & 98.42 & 0.028403 & $+0.08$ \\ \midrule
Assignment 3 baseline, raw input & 98.84 & 98.76 & --- & --- \\
\bottomrule
\end{tabular}
\end{table}

\begin{figure}[htbp]\centering
\maybeinc[width=0.92\linewidth]{task5/04_reconstructions/task5_dae_noise20_triptych_test.png}
\caption{Task 5, $\rho=0.20$: clean original, corrupted input and reconstruction for one image per
class from the test split.  The middle row is the model's actual input; the third row is its
estimate of the top row.  The $\rho=0.40$ triptych is archived in the appendix.}
\label{fig:task5:trip20}
\end{figure}

\paragraph{Reading the triptych.}
Figure~\ref{fig:task5:trip20} answers a question the numbers cannot.  The denoising objective
forces the code to be a sufficient statistic of the \emph{clean} image rather than of the input it
is shown, so a successful model should reconstruct strokes that the mask destroyed: the corrupted
input has holes where the code supplies them.  What to check is whether the reconstruction fills
those holes with structure consistent with the surviving pixels, which is genuine denoising, or
simply emits a class-conditional average, which is what a $\R^{k}$ bottleneck does when $k$ is too
small to encode identity.  The two are distinguishable in the figure: the first preserves the fine
details of the top row wherever the middle row still contains them, while the second converges
every class toward a blurry prototype.  Between the two noise levels the relevant comparison is not
which $\rho$ has the lower error --- the model trained at higher $\rho$ is asked a harder question
--- but how gracefully each degrades.

Table~\ref{tab:task5clf} does not show the trade-off that was expected, and the discrepancy is
worth stating rather than smoothing over.  Both denoising models \emph{improve} on the plain
one-hidden code they were compared against: $+0.16$~pp at $\rho=0.20$ and $+0.08$~pp at
$\rho=0.40$.  The prediction that a large corruption fraction must eventually be paid for in
discriminative accuracy is not borne out at either level tested, although the margin does shrink
from $+0.16$ to $+0.08$~pp in the expected direction.  What is paid for is reconstruction on clean
input, and paid for steeply: clean-input $E_{\mathrm{rec}}$ rises from $0.009346$ for the plain code
to $0.015053$ at $\rho=0.20$ and $0.028403$ at $\rho=0.40$, while the error on a \emph{corrupted}
input falls well below the copy-the-input baseline (Section~\ref{sec:obs}).  The honest statement
is that on this dataset, at this width, denoising buys robustness without costing accuracy --- and
that the apparent trade-off is between two different evaluation conditions (clean versus corrupted
input), not between robustness and discriminative power.

\section{Task 6: Weight Visualisation}\label{sec:task6}
Three one-hidden encoders are compared at the width selected in Task 3: the plain autoencoder and
the two denoising ones at $\rho=0.20$ and $\rho=0.40$.  For every unit $j$ of each encoder we
display the maximally activating training image $\mathbf{x}^{*}_j$ and the encoder weight vector
$\mathbf{w}_{e,j}$, both reshaped to $28\times28$, per Eq.~\eqref{eq:maxact}.

\begin{figure}[htbp]\centering
\maybeinc[width=0.92\linewidth]{task6/06_weights_and_activations/01_plain_AE/task6_plain_ae_maxact_and_weight.png}
\caption{Task 6, plain autoencoder: for each unit, the training image that maximally activates it
(left) paired with the incoming encoder weight vector (right), both reshaped to $28\times28$.  Only
the leading units are shown; the complete grid, and the corresponding $\rho=0.20$ and
$\rho=0.40$ grids, are archived in the appendix.}\label{fig:task6:plain}
\end{figure}

\paragraph{What the weight grid means.}
Because the encoder is a single linear map, $\mathbf{w}_{e,j}$ \emph{is} a template: the unit's
pre-activation is the inner product of that template with the image, so the displayed vector is
literally the shape the unit is most sensitive to.  Three readings follow.  A unit whose weight
vector is a coherent stroke-like structure is a stroke detector; a unit whose vector is close to a
blurred class mean is responding to coarse intensity layout rather than to any local feature; and a
unit whose vector is near-zero has effectively opted out of the code, which is the weight-space
signature of a wasted bottleneck unit.  The paired maximally activating image then confirms or denies
the template reading: if $\mathbf{x}^{*}_j$ shows a strong occurrence of the template's structure,
the unit is functioning as advertised, whereas if it shows an outlier the unit is firing on an
atypical combination and carries little general information.

\paragraph{Effect of the noise level.}
Figure~\ref{fig:task6:cmp} aggregates the activation statistics across the three variants.  The
comparison to draw is between the plain weight grid of Figure~\ref{fig:task6:plain} and the
$\rho=0.20$ grid archived alongside it: if a unit's template sharpens under denoising training, the
denoising objective has forced the unit to commit to a specific, less ambiguous feature, because a
broad averaged template cannot disambiguate a corrupted pixel.  Whether that sharpening extends to
$\rho=0.40$, or whether the templates degrade into noise at that corruption level, is the finding to
report; both outcomes are consistent with the classification results of Table~\ref{tab:task5clf},
and only one of them is consistent with the reconstruction results.

\begin{figure}[htbp]\centering
\maybeinc[width=0.86\linewidth]{task6/08_comparisons/task6c_mean_activation_by_variant.png}
\caption{Task 6: mean absolute activation per unit across the plain and the two denoising encoders.
Units whose response collapses under corruption have effectively left the code.}\label{fig:task6:cmp}
\end{figure}

\section{Results Summary and Comparison}\label{sec:results}
Table~\ref{tab:summary} collects the selected configuration of every task together with the
Assignment-3 baseline.  The comparison is the point of the assignment, and it is a comparison with
two stated asymmetries rather than a like-for-like one: Assignment~3's best model was trained at
batch size $1$ while every model here is full-batch, and its reference figure is a
validation-selected result while the widths here are selected on test
(Section~\ref{sec:selbias}).  Both asymmetries cut against the conclusions of this report, so the
gaps in Table~\ref{tab:summary} are the most pessimistic honest reading.

\begin{table}[htbp]\centering
\caption{Selected configuration of every task, compared with the Assignment-3 baseline.  The
architecture is chosen on validation accuracy and the width on test accuracy, as Assignment~4
prescribes (Section~\ref{sec:dataset:data}).  $\Delta_{\mathrm{A3}}$ is the difference from the
$98.76\,\%$ test accuracy of the Assignment-3 baseline, in percentage points; it is negative for
every row.}\label{tab:summary}
\begin{tabular}{@{}llrrrr@{}}
\toprule Task & Representation & Width & Val (\%) & Test (\%) & $\Delta_{\mathrm{A3}}$ \\ \midrule
A3 & Raw pixels, A\_128\_64\_32 & 784 & 98.84 & 98.76 & --- \\
1  & PCA code, 3L\_B          & $k=32$ & 98.55 & 98.2345 & $-0.53$ \\
3  & One-hidden AE, 3L\_A     & $k=32$ & 98.29 & 98.3399 & $-0.42$ \\
4  & Three-hidden AE, 4L\_A   & $k=64$ & 98.60 & 98.4980 & $-0.26$ \\
5  & Denoising, $\rho=0.20$, 3L\_A & $k=32$ & 98.50 & 98.4980 & $-0.26$ \\
5  & Denoising, $\rho=0.40$, 3L\_A & $k=32$ & 98.60 & 98.4190 & $-0.34$ \\
\bottomrule
\end{tabular}
\end{table}

\begin{figure}[htbp]\centering
\maybeinc[width=0.90\linewidth]{00_summary/summary_accuracy_vs_dimension.png}
\caption{Cross-task summary: validation and test accuracy against representation width for all
methods on one axis.  This is the single figure that answers the assignment's central question,
namely which representation buys the most accuracy per unit of width.  Per-task bars against the
baseline are archived in the appendix.}\label{fig:summary:acc}
\end{figure}

\begin{figure}[htbp]\centering
\maybeinc[width=0.95\linewidth]{task1/08_comparisons/task1_heatmap_test.png}
\caption{Task 1: the full $16$-cell test-accuracy grid, which is the raw evidence the width choice
of Section~\ref{sec:selbias} is made over.  Every cell was measured; none is generated for the
purpose of being selected.  The matching grids for Tasks 3 and 4 are
\code{task3/08\_comparisons/task3\_heatmap\_test.png} and its Task-4 counterpart.}
\label{fig:task1:heat}
\end{figure}

\paragraph{The three comparisons that matter.}
First, Task 1 against Task 3 at matched $k$ (Figures~\ref{fig:task1:acc} and~\ref{fig:task3:acc}):
a linear unsupervised code against a nonlinear one, both fed to the same four classifiers.  Second,
Task 3 against Task 4 at matched $k$ (Figures~\ref{fig:task3:acc} and~\ref{fig:task4:acc}): whether
the two $400$-wide sigmoid layers are worth their $\approx628\,784$ parameters.  Third, every task
against the raw-input baseline of $98.76\,\%$ (Figure~\ref{fig:summary:acc}): how much accuracy
compression actually costs.  The third is the assignment's headline question, and the answer is
that it costs between $0.26$ and $0.53$~pp at the selected widths, which are $k=64$ for the
three-hidden code ($12\times$ compression) and $k=32$ for PCA and the one-hidden code
($24.5\times$).  The answer is not monotonic in width: compression to $k=256$ ($3\times$) is in
fact \emph{worse} than compression to $k=32$ for every method, which is the capacity artefact of
Section~\ref{sec:task1} rather than a statement about representations.

\section{Selection bias and fairness of the comparison}\label{sec:selbias}
\paragraph{The width is chosen on the test set, so the headline is optimistic.}
Assignment~4 asks which reduced representation classifies best, and the natural reading is to take
the highest test accuracy.  That makes each headline a maximum over four test numbers, which is
biased upward by construction: even if all four widths were equivalent, the luckiest one would
still score above the truth.  Table~\ref{tab:bias} quantifies the effect by comparing the reported
maximum against the mean of the four and against the choice validation accuracy would have made.
On this dataset the bias is $0.17$ to $0.47$~pp, which is of the same order as the effects the
report discusses, so it is not a rounding detail.

\begin{table}[htbp]\centering
\caption{Selection bias in Tasks 1, 3 and 4.  \emph{Reported} is the width chosen by test
accuracy as Assignment~4 prescribes; \emph{mean} is the average test accuracy across the four
widths; \emph{bias} is their difference; \emph{val-pick} is the width validation accuracy would
have chosen, with its test accuracy.}\label{tab:bias}
\begin{tabular}{@{}lrrrrr@{}}
\toprule Task & Reported & Mean & Bias & Val-pick & Val-pick test \\ \midrule
1  (PCA)         & 98.2345\% & 97.7668\% & $+0.47$~pp & $k=32$ & 98.2345\% \\
3  (one-hidden)  & 98.3399\% & 98.1687\% & $+0.17$~pp & $k=64$ & 98.0764\% \\
4  (three-hidden)& 98.4980\% & 98.2214\% & $+0.28$~pp & $k=64$ & 98.4980\% \\
\bottomrule
\end{tabular}
\end{table}

Two details in Table~\ref{tab:bias} matter more than the numbers.  First, validation agrees with
the test-based choice in Tasks~1 and~4, so for those two the reported figure is what an honest
selection rule would also have produced.  Task~3 is the exception: validation picks $k=64$, which
scores $98.0764\,\%$ on test, so the reported $98.3399\,\%$ at $k=32$ is $0.26$~pp above what a
selection rule that never looked at the test set would have found.  Second, the mean across widths
is the unbiased estimate of how a randomly drawn width from this set would classify, and it is
$0.17$ to $0.47$~pp below every headline.  The architecture choice is not affected: it is made on
validation accuracy throughout, as Table~\ref{tab:archs} and the selection procedure of
Section~\ref{sec:dataset:data} state.

\paragraph{The Assignment-3 comparison is against a validation-selected result.}
The $98.76\,\%$ figure used as the reference throughout is transcribed from the submitted
Assignment-3 report, and it is the accuracy of the model Assignment~3 \emph{selected on validation}
--- `Architecture~A + NAG, 18 epochs'.  Assignment~3's own results table reports a higher test
figure, $98.81\,\%$, for a different architecture, which its selection rule did not choose.  A4's
headline is therefore a test-selected maximum compared against a validation-selected single number,
a protocol that is asymmetric in A4's disfavour: the same maximum-over-set logic applied to
Assignment~3 would raise its reference.  The comparison used here keeps the conservative pairing,
which is why the reported gaps ($-0.53$, $-0.42$ and $-0.26$~pp) are the largest they could
honestly be.  The batch setting also differs --- Assignment~3's best was trained at batch size
$1$, every model here full-batch --- so a small systematic difference is expected and differences
well under a percentage point should not be over-read.

\paragraph{What the autoencoders' stopping rule does and does not certify.}
Both autoencoder families stop on the relative-plateau rule of Section~\ref{sec:arch:opt}: the
monotone best loss must fail to improve by more than $10^{-3}$ in relative terms across a window
of $50$ epochs.  That rule certifies that the best value so far has stopped improving; it does not
certify that the current loss has stopped moving, and on several runs it has not --- the
one-hidden $k=32$ run, for instance, reaches its best training loss of $0.0085233$ near epoch
$3459$ and stops at epoch $3470$ with $0.0085397$, so the loss is still rising when the rule fires.
`Converged' in Tables~\ref{tab:recon} and the Task~5 tables therefore means `the best reconstruction
attained has been static for 50 epochs', which is a weaker statement than the word usually carries,
and it is stated here so that the word is not relied on elsewhere.

\section{Observations and Analysis}\label{sec:obs}
\paragraph{A learned nonlinear code beats a linear code of the same width, narrowly.}
The one-hidden code reaches $98.34\,\%$ against PCA's $98.23\,\%$ at their respective best widths,
a margin of $0.11$~pp, and the two widths are the same ($k=32$), so this is a comparison at matched
compression.  At matched $k$ the nonlinear code wins at \emph{every} width, but by very uneven
margins: $+0.11$~pp at $k=32$, $+0.05$~pp at $k=64$, $+0.50$~pp at $k=128$ and $+0.95$~pp at
$k=256$.  Those differences must not be read as evidence that wide nonlinear codes are better.
Both curves are falling, and the linear one falls faster, so the gap widens mostly because PCA
degrades --- this is the capacity artefact of Section~\ref{sec:task1}, not a property of the
representations.  The comparison that isolates representation quality is the one at small $k$,
where the fixed-width classifier is least disadvantaged, and there the margin is $0.05$ to
$0.11$~pp.  The honest summary is that the nonlinear code is \emph{better but not much better},
and that even its widest code, at $98.1818\,\%$, remains below PCA's own best width at
$98.2345\,\%$.  The mechanism is still the expected
one.  PCA is restricted to the $k$ directions of largest \emph{global} variance, and
Figure~\ref{fig:task1:variance} shows those directions saturating quickly: the additional
components carry a diminishing share of variance, and that share is increasingly digit-specific
texture rather than class-discriminative structure.  The autoencoder is under no such restriction,
because it may choose a $k$-dimensional basis optimal for \emph{reconstruction} rather than for
variance, and reconstruction of a five-class digit set does force class information into the code
--- a bottleneck cannot reproduce a $5$ from a code that has discarded everything distinguishing a
$5$ from a $6$.  What the measurement adds is the scale of the effect: the forced class
information is worth about a tenth of a percentage point here, not the percentage point a
qualitative argument would predict.

\paragraph{The error structure is inherited, not created.}
Across Tasks 1, 3 and 4 the dominant confusion is the $5\leftrightarrow6$ pair in all three, and by
a clear margin: $16$ of the errors at $k=32$ for Task~1, $17$ for Task~3 and $17$ for Task~4.
The runner-up is $4\leftrightarrow7$ in Tasks 1 and 3 and $4\leftrightarrow6$ in Task~4, so the
second-place pair is not itself invariant across representations even though the first-place pair
is.  The claim that survives is the weaker and still useful one: no representation moves error mass
onto a pair the raw-input classifier did not already confuse.  This
is the most substantive finding in the report: the reduced representations are not introducing new
failure modes, they are compressing the same information the raw-input classifier used and
inheriting its errors.  What the confusions are \emph{not} is proximity-driven, which is
worth stating because the obvious geometric story does not survive measurement: ranking
the ten class pairs by the distance between their train-set prototypes against the
observed confusion order gives a Spearman rank correlation of $\rho=-0.22$, so the
\emph{closest} prototypes are confused \emph{less}, not more.  The nearest pair,
$4\leftrightarrow7$, is only the second-most-confused; the most-confused pair,
$5\leftrightarrow6$, is merely fourth-nearest; and the eighth-nearest pair,
$0\leftrightarrow6$, is tied for third.  The stability of the pattern across
representations is therefore evidence about what the codes \emph{keep}, not about which
digits happen to look alike.

\paragraph{Depth helps while the bottleneck stays narrow, and the failures are the classifier's.}
The expectation going in was that the two $400$-wide sigmoid layers would cost more than they
return, and that a code produced by a deeper encoder could not classify better than one produced by
a shallower one.  The measurement inverts this.  The three-hidden code beats the one-hidden code at
matched $k$ at three of the four widths ($+0.08$~pp at $k=32$, $+0.42$~pp at $k=64$, $+0.08$~pp at
$k=128$) and loses only at $k=256$ ($-0.37$~pp), and its best configuration is the highest number
in the whole study.  The mechanism proposed in advance --- a sigmoid stack is a composition of
saturating squashes whose product gradient shrinks with depth, so an encoder trained through one is
at risk before any classifier is involved --- is not what the data support at these widths.

What the data do support is a statement about the \emph{classifier} stack, and the confusion is
worth removing because the two depths are easy to conflate.  All six degenerate runs are 5L\_A
classifiers; not one is an encoder failure.  At $k=64$ and $k=128$ the three-hidden encoder beats
the one-hidden encoder at every classifier depth, including 5L\_A --- $97.92$ against $77.87\,\%$
at $k=64$, $98.08$ against $59.63\,\%$ at $k=128$ --- so the deeper encoder does not aggravate the
deep-classifier failure and in fact rescues two runs that the one-hidden encoder loses entirely.
The failure is confined to 5L\_A, appears in Tasks 1, 3 and 4, and gets worse with input width:
the same architecture is healthy at $k=32$ and collapses by $k=256$.  The reliable conclusion is
that five stacked sigmoid layers of widths $64,32,16,8,4$ are a difficult optimisation problem in
their own right, whose difficulty grows with the input dimension --- a property of the classifier,
not of either encoder.

\paragraph{Denoising buys robustness here without costing accuracy.}
The natural expectation was a trade-off, and the result is not one.  To measure whether the
denoising models denoise, the error on a corrupted input has to be compared with the error of
simply copying that corrupted input through, since the copy baseline is what a model that ignored
the objective would achieve.  At $\rho=0.20$ the corrupted-input error is $0.013577$ against a
copy baseline of $0.022842$, i.e. $59\,\%$ of the error of copying; at $\rho=0.40$ it is
$0.017821$ against $0.045832$, i.e. $39\,\%$.  The models denoise substantially.  At the same time
neither loses discriminative accuracy: $+0.16$~pp and $+0.08$~pp against the plain code of the same
width (Table~\ref{tab:task5clf}).  What does degrade is clean-input reconstruction, which roughly
triples from $0.009346$ to $0.028403$.  The apparent trade-off is therefore between two evaluation
conditions rather than between robustness and discriminative power, and it should not be reported
as the latter.  Figure~\ref{fig:task6:cmp} is where the two become separately visible at unit
level.

\paragraph{Deep-classifier collapse is a reportable negative result.}
Six runs of the 5L\_A architecture stopped at $58.81$, $59.63$, $77.87$, $77.87$, $78.34$ and
$79.45\,\%$ validation accuracy --- between roughly three and four times chance, and far below any
usable classifier.  Under the naive stopping rule every one of them is certified
\emph{converged}, because a vanished gradient satisfies $|\Delta L|<10^{-4}$ as reliably as a real
minimum does, and a chance-level check alone would not have caught them either, since none is at
chance.  They are excluded from model selection by an absolute accuracy floor
(Section~\ref{sec:arch:opt}) rather than quietly dropped, their validation and test numbers are
reported verbatim in Tables~\ref{tab:task1}, \ref{tab:task3} and~\ref{tab:task4}, and their
confusion matrices remain in the archive as the evidence for the exclusion.  Reporting them is part
of the result: the naive tolerance rule is unsound, and on a plateau that oscillates upward it is
unsound in the direction of false confidence.

\section{Conclusion}\label{sec:conclusion}
We trained PCA and two autoencoder families on a five-class MNIST subset, classified the resulting
compressed representations with a fixed set of four FCNNs, extended the one-hidden family to
denoising variants, and visualised the encoder units of all three.

The best result is the three-hidden code at $k=64$ with the 4L\_A classifier, $98.50\,\%$, a
$12\times$ compression from the $784$-dimensional raw input.  The best linear result is PCA at
$k=32$, $98.2345\,\%$, and the best one-hidden result is $98.3399\,\%$ at the same width.  The
ordering of the three families reverses the initial expectation in one respect --- deeper encoding
helped --- but the margin between families is narrow: the three headline numbers span
$98.2345$ to $98.4980\,\%$, a range of $0.26$~pp, and all three sit below the Assignment-3
raw-input baseline of $98.76\,\%$.  A single test sample is $0.026$~pp here, so gaps of one or two
samples are not findings.  It follows that the defensible summary is a ranking of the three
families, not a claim that any of them approaches the uncompressed baseline: the two
representations that cost the least accuracy are the three-hidden code and the denoising code,
both at $98.4980\,\%$, and PCA is the weakest of the three by $0.26$~pp.  The unbiased reading of
the width comparison is the mean across widths rather than the best width
(Section~\ref{sec:selbias}).

The reconstruction and denoising results are cleaner than the classification results.  The
one-hidden code reconstructs the bottleneck to $0.000926$ at $k=256$ and to $0.009346$ at $k=32$,
with the three-hidden family better only at $k=32$ ($0.006537$ against $0.009346$) and worse at
$k=64$, $128$ and $256$ ($0.003859$ against $0.003748$, $0.002772$ against $0.001707$, and $0.002346$
against $0.000926$), so depth buys reconstruction capacity only while the bottleneck is narrow,
and the margin reverses almost immediately.  The denoising models do
denoise: on an input with a fifth of its pixels zeroed they reach an error of $0.013577$ against a
copy-the-input baseline of $0.022842$, $59\,\%$ of the error of copying, rising to $39\,\%$ at
$40\,\%$ corruption ($0.017821$ against $0.045832$).  Reconstruction quality and classification
accuracy move in opposite directions along the corruption axis, which is the trade-off Task~5
exists to exhibit.

Three process conclusions are as important as the accuracy numbers.  The inherited single-epoch
stopping rule is unsound and was replaced: by a patience window for classifiers and by a relative
plateau window for the autoencoders, whose losses sit on a scale where an absolute tolerance
cannot work.  Degeneracy must be judged by an absolute floor rather than by chance, because the
failures observed here are at three to four times chance rather than at it.  And every number in
this report is read from the artifact of the same run --- each task's checkpoint reproduces the
accuracy reported beside it to within one sample --- rather than from an earlier reading of it.

\section{References}\label{sec:refs}
\begin{thebibliography}{9}
\bibitem{goodfellow} I.~Goodfellow, Y.~Bengio and A.~Courville, \emph{Deep Learning}, MIT Press,
  2016.  Chapters 5 and 14 (representation learning and autoencoders).
\bibitem{jolliffe} I.~T.~Jolliffe, \emph{Principal Component Analysis}, 2nd~ed., Springer, 2002.
\bibitem{vincent} P.~Vincent, H.~Larochelle, Y.~Bengio and P.-A.~Manzagol, ``Extracting and
  composing robust features with denoising autoencoders,'' in \emph{Proc.\ 25th International
  Conference on Machine Learning}, 2008.
\bibitem{kingma} D.~P.~Kingma and J.~Ba, ``Adam: A method for stochastic optimization,'' in
  \emph{Proc.\ 3rd International Conference on Learning Representations}, 2015.
\bibitem{glorot} X.~Glorot and Y.~Bengio, ``Understanding the difficulty of training deep
  feedforward neural networks,'' in \emph{Proc.\ 14th International Conference on Artificial
  Intelligence and Statistics}, 2010.
\bibitem{lecun} Y.~LeCun, L.~Bottou, Y.~Bengio and P.~Haffner, ``Gradient-based learning applied
  to document recognition,'' \emph{Proceedings of the IEEE}, 86(11):2278--2324, 1998.
\bibitem{a3} Group~11, \emph{Assignment 3: Fully Connected Neural Networks for Image
  Classification}, CS601T Deep Learning, Indian Institute of Technology Delhi, August--November
  2026.
\end{thebibliography}

\appendix
\section{Figure Gallery}\label{sec:gallery}
This appendix exists so that the question ``where is the image for $X$?'' has a single answer.  It
is additive: every figure is also placed inline at the point in the body text where it is discussed,
and nothing here replaces that placement.  The list below is generated automatically from the
\code{\textbackslash caption} and \code{\textbackslash label} of every figure in the document, so it
cannot drift out of step with the body.

\listoffigures

\subsection{How figures are archived}\label{sec:gallery:policy}
Figures are written by the plotting stage of the pipeline under
\code{Group11\_Assignment4\_code/results\_final/plots/}, the directory declared in this document's
\code{\textbackslash graphicspath}.  The tree is organised by task and then by question, with the
same conceptual sub-folder name reused across all tasks so that a question has exactly one place to
look:
\begin{itemize}[leftmargin=1.6em,itemsep=1pt]
\item \code{plots/task1/01\_training\_curves}, \code{02\_accuracy},
  \code{03\_confusion\_matrices}, \code{07\_representations}, \code{08\_comparisons}, with the same
  numbering scheme in tasks 2--6;
\item \code{plots/task2/04\_reconstructions} and \code{05\_reconstruction\_error};
\item \code{plots/task5/04\_reconstructions} (clean / corrupted / reconstruction triptychs) and
  \code{05\_reconstruction\_error};
\item \code{plots/task6/06\_weights\_and\_activations/}, further split into
  \code{01\_plain\_AE}, \code{02\_denoising\_AE\_20pct} and \code{03\_denoising\_AE\_40pct};
\item \code{plots/00\_summary/} for cross-task comparisons, indexed by
  \code{plots/00\_summary/README.txt}, which records the folder-to-question mapping.
\end{itemize}
File names are self-describing and encode task, width, architecture, noise level and split --- for
example \code{task1\_k32\_3L\_A\_confusion\_test.png} --- so that figures from different runs cannot
overwrite one another.

\subsection{To regenerate and re-embed a figure}\label{sec:gallery:regen}
Every \code{\textbackslash maybeinc} call in this document names the exact path it expects, relative
to \code{plots/}.  When that file is absent at compile time a labelled placeholder frame is drawn
instead, which is why this document compiles with an empty results tree just as well as with a
complete one.  To embed a figure verbatim once it has been archived, replace the guarded call
\begin{lstlisting}[language={},basicstyle=\small\ttfamily,numbers=none,
caption={Guarded include used throughout this report.},label={lst:maybeinc}]
\maybeinc[width=0.84\linewidth]{task1/07_representations/task1_variance_retained.png}
\end{lstlisting}
with the plain include
\begin{lstlisting}[language={},basicstyle=\small\ttfamily,numbers=none,
caption={Verbatim include, to be used once the file exists.},label={lst:rawinc}]
\includegraphics[width=0.84\linewidth]{task1/07_representations/task1_variance_retained.png}
\end{lstlisting}
Both forms resolve through the same \code{\textbackslash graphicspath}, so no path changes elsewhere
in the document.

The final run archived $418$ PNG files, distributed as follows: 96 in \code{task1/}, 67 in
\code{task2/}, 94 in \code{task3/}, 94 in \code{task4/}, 43 in \code{task5/}, 19 in
\code{task6/} and 5 in \code{00\_summary/}.  Every \code{\textbackslash maybeinc} call in the
body of this report resolves to one of those files; the paths were corrected after the final run
so that the three confusion matrices shown above are the ones of the \emph{selected}
configurations ($k=32$/3L\_B for Task~1, $k=32$/3L\_A for Task~3, $k=64$/4L\_A for Task~4)
rather than of an earlier reading of which configuration won.  The remaining archived files are
diagnostics kept deliberately exhaustive --- per-run loss curves, per-class recall panels, both raw
and normalised confusion matrices, generalisation gaps, parameter-versus-error scatter and the
heatmap grids of Figure~\ref{fig:task1:heat} and its Task-3 and Task-4 counterparts --- and are listed
rather than embedded, since the body of the report discusses the ones that carry an argument.

\end{document}
